\chapter{Maxwell, Yang--Mills and Fluids}
\label{ch:fields}

\begin{usesbox}
\textbf{This chapter uses:} quaternions and $\su(2)$ (\cref{ch:quaternions}); the mass shell and the hemisphere point $(m/E,\mathbf p/E)$ (\cref{thm:hemibij}); connections and their curvature (\cref{ch:berrywigner,ch:monopoles}); the coordinate formulas and the two gauges of \cref{ch:sphmetrics}.\\
\textbf{It provides:} Maxwell's equations as the Euler--Lagrange equations of the Maxwell action plus the Bianchi identity; charged particles and dust; relativistic perfect fluids, Kelvin's theorem and flux freezing; Yang--Mills fields, their gauge covariance and first variation; the three homogeneous Yang--Mills fields on $S^3$; and the field strength and stress tensor of the spherically symmetric magnetic Yang--Mills ansatz in both gauges.
\end{usesbox}

This chapter assembles the matter content of the rest of the book. Each field is introduced through its action, and its field equations are obtained as Euler--Lagrange equations. Each conservation law is then derived either as an identity that holds off shell, such as the Bianchi identity, or as a consequence of the field equations. Keeping these two kinds of statement apart is the point of the presentation, and the formal verification enforces it: in every theorem below it is stated explicitly which equations are assumed.

Throughout, $\eta=\operatorname{diag}(-1,1,1,1)$, derivatives commute, and the coupling constants are constant.

\section{Maxwell's equations}

The field strength of a potential $A_\mu$ is $F_{\mu\nu}=\partial_\mu A_\nu-\partial_\nu A_\mu$. Half of Maxwell's equations hold for every potential.

\begin{theorem}[The Bianchi identity and its consequences]\label{thm:bianchi}\ver\
For every $A_\mu$: $\partial_\lambda F_{\mu\nu}+\partial_\mu F_{\nu\lambda}+\partial_\nu F_{\lambda\mu}=0$. In terms of $\mathbf E$ and $\mathbf B$ this is $\nabla\cdot\mathbf B=0$ and $\nabla\times\mathbf E+\partial_t\mathbf B=0$. Moreover $\partial_\nu\partial_\mu F^{\mu\nu}=0$ identically.
\end{theorem}
\begin{proof}
Substitute $F_{\mu\nu}=\partial_\mu A_\nu-\partial_\nu A_\mu$ into the cyclic sum:
\[
\begin{aligned}
&\partial_\lambda\partial_\mu A_\nu-\partial_\lambda\partial_\nu A_\mu
+\partial_\mu\partial_\nu A_\lambda-\partial_\mu\partial_\lambda A_\nu
+\partial_\nu\partial_\lambda A_\mu-\partial_\nu\partial_\mu A_\lambda .
\end{aligned}
\]
Since partial derivatives commute, the first and fourth terms cancel, the second and fifth cancel, and the third and sixth cancel. So the cyclic sum is zero.

To read off the vector form we use the conventions forced by the components of $\mathcal E$ in \cref{thm:maxwellvar}: $\eta=\operatorname{diag}(-1,1,1,1)$, $E_i=F_{i0}$ and $B_i=\frac12\epsilon_{ijk}F_{jk}$, so that $F_{12}=B_3$, $F_{23}=B_1$, $F_{31}=B_2$.
\begin{itemize}[nosep]
\item $(\lambda,\mu,\nu)=(1,2,3)$: $\partial_1F_{23}+\partial_2F_{31}+\partial_3F_{12}=\partial_1B_1+\partial_2B_2+\partial_3B_3=\nabla\cdot\mathbf B$.
\item $(\lambda,\mu,\nu)=(0,2,3)$: $\partial_0F_{23}+\partial_2F_{30}+\partial_3F_{02}=\partial_tB_1+\partial_2E_3-\partial_3E_2=\partial_tB_1+(\nabla\times\mathbf E)_1$.
\end{itemize}
The triples $(0,3,1)$ and $(0,1,2)$ give the second and third components in the same way.

For the last claim, write $S=\partial_\nu\partial_\mu F^{\mu\nu}$. Renaming the summation indices $\mu\leftrightarrow\nu$ gives $S=\partial_\mu\partial_\nu F^{\nu\mu}$. Using $F^{\nu\mu}=-F^{\mu\nu}$ and the commutation of partial derivatives, this equals $-\partial_\nu\partial_\mu F^{\mu\nu}=-S$. Hence $S=0$.
\end{proof}


The other half comes from the action $S=\int\bigl(-\frac14F_{\mu\nu}F^{\mu\nu}+J^\mu A_\mu\bigr)$.

\begin{theorem}[First variation of the Maxwell action]\label{thm:maxwellvar}\ver\
Under $A\mapsto A+\varepsilon B$ with constant $\varepsilon$, the Lagrangian changes by $\varepsilon$ times $\sum_\nu\mathcal E^\nu B_\nu$ plus a total divergence, plus a term of order $\varepsilon^2$. The Euler--Lagrange expression $\mathcal E^\nu$ has components
\[
\mathcal E^0=-\nabla\cdot\mathbf E+\rho,\qquad \mathcal E^1=\partial_tE_1-(\nabla\times\mathbf B)_1+J^1,
\]
so $\mathcal E=0$ is Gauss's law $\nabla\cdot\mathbf E=\rho$ and the Amp\`ere--Maxwell law $\nabla\times\mathbf B-\partial_t\mathbf E=\mathbf J$. If $\mathcal E=0$ then $\partial_\nu J^\nu=0$.
\end{theorem}
\begin{proof}
Write $F_B=\partial_\mu B_\nu-\partial_\nu B_\mu$. Then $F(A+\varepsilon B)=F+\varepsilon F_B$, and
\[
-\tfrac14(F+\varepsilon F_B)^2=-\tfrac14F^2-\tfrac12\varepsilon\,F^{\mu\nu}(F_B)_{\mu\nu}-\tfrac14\varepsilon^2F_B^2 .
\]
The $\varepsilon^2$ term is the one announced. For the linear term, the antisymmetry of $F^{\mu\nu}$ gives
\[
F^{\mu\nu}(\partial_\mu B_\nu-\partial_\nu B_\mu)=2F^{\mu\nu}\partial_\mu B_\nu .
\]
Hence
\[
-\tfrac12F^{\mu\nu}(F_B)_{\mu\nu}=-F^{\mu\nu}\partial_\mu B_\nu=(\partial_\mu F^{\mu\nu})B_\nu-\partial_\mu\bigl(F^{\mu\nu}B_\nu\bigr).
\]
Adding the source term $\varepsilon J^\nu B_\nu$, the coefficient of $\varepsilon B_\nu$ is
\[
\mathcal E^\nu=\partial_\mu F^{\mu\nu}+J^\nu ,
\]
and the total divergence is $-\partial_\mu(F^{\mu\nu}B_\nu)$.

Now take components. Raising indices with $\eta=\operatorname{diag}(-1,1,1,1)$ gives $F^{i0}=-F_{i0}=-E_i$ and $F^{ij}=F_{ij}$.
\begin{itemize}[nosep]
\item $\nu=0$: $\mathcal E^0=\partial_iF^{i0}+J^0=-\nabla\cdot\mathbf E+\rho$.
\item $\nu=1$: $\mathcal E^1=\partial_0F^{01}+\partial_2F^{21}+\partial_3F^{31}+J^1$. Here $F^{01}=-F_{01}=F_{10}=E_1$, $F^{21}=F_{21}=-B_3$ and $F^{31}=F_{31}=B_2$. So
\[
\mathcal E^1=\partial_tE_1-(\partial_2B_3-\partial_3B_2)+J^1=\partial_tE_1-(\nabla\times\mathbf B)_1+J^1 .
\]
\end{itemize}
Setting $\mathcal E=0$ gives Gauss's law and the Amp\`ere--Maxwell law.

Finally, take the divergence of $\mathcal E^\nu=\partial_\mu F^{\mu\nu}+J^\nu$. The first term is $\partial_\nu\partial_\mu F^{\mu\nu}=0$ by \cref{thm:bianchi}, so $\partial_\nu\mathcal E^\nu=\partial_\nu J^\nu$. If $\mathcal E=0$, then $\partial_\nu J^\nu=0$.
\end{proof}


Charge conservation is not an extra postulate: it is the compatibility condition between the field equations and the identity $\partial_\nu\partial_\mu F^{\mu\nu}=0$. In spherical symmetry Gauss's law integrates to the familiar statement that the flux through a sphere measures the enclosed charge.

\begin{proposition}[Integrated Gauss law]\label{prop:gaussint}\ver\
If $\partial_r(r^2E_r)=r^2\rho$ on $[a,b]$, then $4\pi b^2E_r(b)-4\pi a^2E_r(a)=4\pi\int_a^br^2\rho\,\dd r$.
\end{proposition}
\begin{proof}
The function $r\mapsto r^2E_r(r)$ has derivative $r^2\rho$ on $[a,b]$. By the fundamental theorem of calculus,
\[
b^2E_r(b)-a^2E_r(a)=\int_a^br^2\rho\,\dd r .
\]
Multiplying by $4\pi$, the area of the unit sphere, gives the claim. The left side is the outward flux through the sphere of radius $b$ minus the flux through the sphere of radius $a$.
\end{proof}


The Maxwell stress tensor $T_{\mu\nu}=F_{\mu\alpha}F_\nu{}^\alpha-\frac14\eta_{\mu\nu}F^2$ satisfies the identity $\partial^\mu T_{\mu\nu}=F_{\nu\alpha}\,\partial_\mu F^{\mu\alpha}$, which holds once the Bianchi identity is used; on shell the right side is $-F_{\nu\alpha}J^\alpha$, the Lorentz force density.

\section{Charged particles and dust}

\begin{theorem}[Charged particle]\label{thm:particle}\ver\
For the worldline Lagrangian $L=\frac12m\,\eta_{\mu\nu}u^\mu u^\nu+q\,A_\mu u^\mu$, with $u^\mu=\dd x^\mu/\dd\tau$ and the chain rule $\dd A_\mu/\dd\tau=u^\nu\partial_\nu A_\mu$:
\begin{enumerate}[label=(\roman*)]
\item The Euler--Lagrange equations are the Lorentz force law $m\,\dd u_\mu/\dd\tau=qF_{\mu\nu}u^\nu$.
\item With $p_\mu=m\eta_{\mu\nu}u^\nu+qA_\mu$ and $H=\frac1{2m}\eta^{\mu\nu}(p_\mu-qA_\mu)(p_\nu-qA_\nu)$, Hamilton's equations are equivalent to (i).
\item If $\eta_{\mu\nu}u^\mu u^\nu=-1$ with $u^0>0$, then $E_{\mathrm{kin}}=mu^0$ and $\mathbf p_{\mathrm{kin}}=m\mathbf u$ satisfy $E_{\mathrm{kin}}^2=m^2+|\mathbf p_{\mathrm{kin}}|^2$, and $(m/E_{\mathrm{kin}},\mathbf p_{\mathrm{kin}}/E_{\mathrm{kin}})$ is a point of the unit sphere.
\end{enumerate}
\end{theorem}
\begin{proof}
\emph{(i)} Treat $x^\mu$ and $u^\mu$ as independent in the Lagrangian.
\begin{itemize}[nosep]
\item Velocity derivative: $\partial L/\partial u^\mu=m\,u_\mu+qA_\mu$. Its $\tau$-derivative, by the chain rule, is
\[
\frac{\dd}{\dd\tau}\frac{\partial L}{\partial u^\mu}=m\frac{\dd u_\mu}{\dd\tau}+q\,u^\nu\partial_\nu A_\mu .
\]
\item Position derivative: $A$ is the only place $x$ appears, so $\partial L/\partial x^\mu=q\,u^\nu\partial_\mu A_\nu$.
\end{itemize}
The Euler--Lagrange equation sets the two equal:
\[
m\frac{\dd u_\mu}{\dd\tau}=q\,u^\nu(\partial_\mu A_\nu-\partial_\nu A_\mu)=qF_{\mu\nu}u^\nu .
\]

\emph{(ii)} Put $\pi_\mu=p_\mu-qA_\mu$, so that $H=\frac1{2m}\eta^{\mu\nu}\pi_\mu\pi_\nu$. Hamilton's equations are
\[
\frac{\dd x^\mu}{\dd\tau}=\frac{\partial H}{\partial p_\mu}=\frac1m\eta^{\mu\nu}\pi_\nu,\qquad
\frac{\dd p_\mu}{\dd\tau}=-\frac{\partial H}{\partial x^\mu}=\frac qm\eta^{\alpha\beta}\pi_\alpha\,\partial_\mu A_\beta .
\]
The first equation says $\pi_\mu=m\,u_\mu$, that is, $p_\mu=m u_\mu+qA_\mu$; this is the definition of $p$ used in (ii). With it, the second equation reads
\[
m\frac{\dd u_\mu}{\dd\tau}+q\,u^\nu\partial_\nu A_\mu=q\,u^\beta\partial_\mu A_\beta ,
\]
which is the equation of (i). Conversely, given a solution of (i), define $p_\mu=mu_\mu+qA_\mu$; the same lines read backwards show that Hamilton's equations hold.

\emph{(iii)} Write $\eta_{\mu\nu}u^\mu u^\nu=-1$ as $-(u^0)^2+|\mathbf u|^2=-1$. Multiplying by $-m^2$ gives
\[
(mu^0)^2=m^2+m^2|\mathbf u|^2,\qquad\text{that is,}\qquad E_{\mathrm{kin}}^2=m^2+|\mathbf p_{\mathrm{kin}}|^2 .
\]
Divide by $E_{\mathrm{kin}}^2>0$: $(m/E_{\mathrm{kin}})^2+|\mathbf p_{\mathrm{kin}}/E_{\mathrm{kin}}|^2=1$. This is a point of the unit sphere in $\RR^4$, and its first coordinate $m/E_{\mathrm{kin}}$ is positive.
\end{proof}


Item (iii) connects the classical particle with \cref{ch:hemisphere}: the kinetic four-momentum of a charged particle in any field, divided by its energy, is a point of the hemisphere.

\begin{proposition}[Charged dust]\label{prop:dust}\ver\
For dust $T^{\mu\nu}=\rho u^\mu u^\nu$ with charge current $\rho_eu^\mu$: if $\partial_\mu(\rho u^\mu)=0$ and the total stress--energy is conserved, then $\rho\,u^\mu\partial_\mu u^\nu=\rho_eF^\nu{}_\lambda u^\lambda$.
\end{proposition}
\begin{proof}
The proof has three steps.

\emph{Step 1: the dust part.} By the product rule,
\[
\partial_\mu(\rho u^\mu u^\nu)=\partial_\mu(\rho u^\mu)\,u^\nu+\rho\,u^\mu\partial_\mu u^\nu=\rho\,u^\mu\partial_\mu u^\nu ,
\]
where we used $\partial_\mu(\rho u^\mu)=0$.

\emph{Step 2: the field part.} We show $\partial^\mu T_{\mu\nu}=F_{\nu\alpha}\partial_\mu F^{\mu\alpha}$ for the Maxwell tensor. The derivative of the first term is
\[
\partial^\mu(F_{\mu\alpha}F_\nu{}^\alpha)=(\partial^\mu F_{\mu\alpha})F_\nu{}^\alpha+F^{\mu\alpha}\partial_\mu F_{\nu\alpha}.
\]
The derivative of the second term is $\partial_\nu(\frac14F^2)=\frac12F^{\mu\alpha}\partial_\nu F_{\mu\alpha}$. The Bianchi identity gives
\[
\partial_\nu F_{\mu\alpha}=-\partial_\mu F_{\alpha\nu}-\partial_\alpha F_{\nu\mu}=\partial_\mu F_{\nu\alpha}-\partial_\alpha F_{\nu\mu}.
\]
Contract with $F^{\mu\alpha}$. In the second term rename $\mu\leftrightarrow\alpha$ and use $F^{\alpha\mu}=-F^{\mu\alpha}$; it becomes equal to the first. Hence
\[
\tfrac12F^{\mu\alpha}\partial_\nu F_{\mu\alpha}=F^{\mu\alpha}\partial_\mu F_{\nu\alpha},
\]
which cancels the second term of the first derivative. What remains is $(\partial^\mu F_{\mu\alpha})F_\nu{}^\alpha=F_{\nu\alpha}\partial_\mu F^{\mu\alpha}$.

\emph{Step 3: combine.} On shell, \cref{thm:maxwellvar} gives $\partial_\mu F^{\mu\alpha}=-J^\alpha$ with $J^\alpha=\rho_eu^\alpha$, so the field part is $\partial_\mu T_{\mathrm{em}}^{\mu\nu}=-\rho_eF^\nu{}_\lambda u^\lambda$. Conservation of the total, $\partial_\mu(T_{\mathrm{dust}}^{\mu\nu}+T_{\mathrm{em}}^{\mu\nu})=0$, together with Step 1 gives $\rho\,u^\mu\partial_\mu u^\nu=\rho_eF^\nu{}_\lambda u^\lambda$.
\end{proof}


\section{Relativistic fluids}

A perfect fluid has $T^{\mu\nu}=(\rho+p)u^\mu u^\nu+p\,\eta^{\mu\nu}$ with $u\cdot u=-1$.

\begin{theorem}[Relativistic Euler equations]\label{thm:euler}\ver\
\begin{enumerate}[label=(\roman*)]
\item $\partial_\mu T^{\mu\nu}=\bigl(\partial_\mu((\rho+p)u^\mu)\bigr)u^\nu+(\rho+p)u^\mu\partial_\mu u^\nu+\eta^{\nu\nu}\partial_\nu p$, and $u\cdot u=-1$ implies $u_\nu\partial_\mu u^\nu=0$.
\item If $\partial_\mu T^{\mu\nu}=0$ then $\partial_\mu((\rho+p)u^\mu)=u^\mu\partial_\mu p$ and
\[
(\rho+p)\,u^\mu\partial_\mu u^\nu=-\eta^{\nu\nu}\partial_\nu p-u^\nu u^\mu\partial_\mu p .
\]
\end{enumerate}
\end{theorem}
\begin{proof}
\emph{(i)} Apply the product rule to $(\rho+p)u^\mu\cdot u^\nu$:
\[
\partial_\mu\bigl((\rho+p)u^\mu u^\nu\bigr)=\partial_\mu\bigl((\rho+p)u^\mu\bigr)u^\nu+(\rho+p)u^\mu\partial_\mu u^\nu .
\]
The metric is constant and diagonal, so $\partial_\mu(p\,\eta^{\mu\nu})=\eta^{\nu\nu}\partial_\nu p$ (no sum over $\nu$). Adding the two gives the formula. Differentiating $u_\nu u^\nu=-1$ gives $2u_\nu\partial_\mu u^\nu=0$.

\emph{(ii)} Contract $\partial_\mu T^{\mu\nu}=0$ with $u_\nu$ and use three facts: $u_\nu u^\nu=-1$, $u_\nu\partial_\mu u^\nu=0$, and $u_\nu\eta^{\nu\nu}\partial_\nu p=u^\nu\partial_\nu p$. The result is
\[
-\partial_\mu\bigl((\rho+p)u^\mu\bigr)+0+u^\nu\partial_\nu p=0 .
\]
This is the first claim. Put it back into $\partial_\mu T^{\mu\nu}=0$:
\[
(u^\mu\partial_\mu p)\,u^\nu+(\rho+p)u^\mu\partial_\mu u^\nu+\eta^{\nu\nu}\partial_\nu p=0 ,
\]
which rearranges to the stated equation.
\end{proof}


Vorticity is governed by Cartan's formula for the Lie derivative. For a closed two-form $\omega=\dd w$, $\mathcal L_u\omega=\dd(i_u\omega)$.

\begin{theorem}[Kelvin and flux freezing]\label{thm:kelvin}\ver\
For $\omega_{\mu\nu}=\partial_\mu w_\nu-\partial_\nu w_\mu$:
\begin{enumerate}[label=(\roman*)]
\item $\partial_\lambda\omega_{\mu\nu}-\partial_\mu\omega_{\lambda\nu}+\partial_\nu\omega_{\lambda\mu}=0$, and $(\mathcal L_u\omega)_{\mu\nu}=\partial_\mu(i_u\omega)_\nu-\partial_\nu(i_u\omega)_\mu$.
\item (Kelvin) If $i_u\omega=0$ then $\mathcal L_u\omega=0$: vorticity is carried along the flow.
\item (Flux freezing) For $F=\dd A$ and ideal magnetohydrodynamics $i_uF=0$, $\mathcal L_uF=0$.
\end{enumerate}
\end{theorem}
\begin{proof}
\emph{(i)} Substitute $\omega_{\mu\nu}=\partial_\mu w_\nu-\partial_\nu w_\mu$:
\[
\partial_\lambda\partial_\mu w_\nu-\partial_\lambda\partial_\nu w_\mu-\partial_\mu\partial_\lambda w_\nu+\partial_\mu\partial_\nu w_\lambda+\partial_\nu\partial_\lambda w_\mu-\partial_\nu\partial_\mu w_\lambda=0 ,
\]
because the six terms cancel in pairs. Equivalently, $\partial_\lambda\omega_{\mu\nu}=\partial_\mu\omega_{\lambda\nu}-\partial_\nu\omega_{\lambda\mu}$.

For the Lie derivative we use the coordinate formula
\[
(\mathcal L_u\omega)_{\mu\nu}=u^\lambda\partial_\lambda\omega_{\mu\nu}+\omega_{\lambda\nu}\partial_\mu u^\lambda+\omega_{\mu\lambda}\partial_\nu u^\lambda .
\]
Since $(i_u\omega)_\nu=u^\lambda\omega_{\lambda\nu}$, the product rule gives
\[
\partial_\mu(i_u\omega)_\nu-\partial_\nu(i_u\omega)_\mu
=\omega_{\lambda\nu}\partial_\mu u^\lambda-\omega_{\lambda\mu}\partial_\nu u^\lambda+u^\lambda\bigl(\partial_\mu\omega_{\lambda\nu}-\partial_\nu\omega_{\lambda\mu}\bigr).
\]
In the second term, $-\omega_{\lambda\mu}=\omega_{\mu\lambda}$. By the first identity, the bracket is $\partial_\lambda\omega_{\mu\nu}$. The right side is therefore the coordinate formula for $(\mathcal L_u\omega)_{\mu\nu}$.

\emph{(ii)} If $i_u\omega=0$, the right side of (i) vanishes.

\emph{(iii)} This is (ii) with $w=A$ and $\omega=F$. The ideal MHD condition $i_uF=0$ says that the electric field vanishes in the rest frame of the fluid.
\end{proof}


In three dimensions the vortex-stretching term involves a cross product, and here quaternions reappear: for pure quaternions $v,w$, $[v,w]=2\,v\times w$. Finally, magnetic helicity obeys the local balance law $\partial_t(\mathbf A\cdot\mathbf B)+\nabla\cdot(\varphi\mathbf B+\mathbf E\times\mathbf A)+2\,\mathbf E\cdot\mathbf B=0$, so it is conserved when $\mathbf E\cdot\mathbf B=0$.

\section{Yang--Mills fields}

For $\SU(2)$ we write the potential as $A^a_\mu$ with $a=1,2,3$ and use the structure constants $\epsilon_{abc}$:
\[
F^a_{\mu\nu}=\partial_\mu A^a_\nu-\partial_\nu A^a_\mu+\epsilon_{abc}A^b_\mu A^c_\nu .
\]
In matrix notation $F=\dd A+A\wedge A$, and a gauge transformation acts by $A\mapsto gAg^{-1}+g\,\dd g^{-1}$.

\begin{theorem}[Gauge covariance and first variation]\label{thm:ymvar}\ver\
\begin{enumerate}[label=(\roman*)]
\item $\partial_k(g^{-1})=-g^{-1}(\partial_kg)g^{-1}$, and under a gauge transformation $F'_{ij}=g\,F_{ij}\,g^{-1}$, in any associative algebra.
\item $F^a_{\nu\mu}=-F^a_{\mu\nu}$, and $\delta F^a_{\mu\nu}=(\nabla_\mu\delta A)^a_\nu-(\nabla_\nu\delta A)^a_\mu$.
\item The Lagrangian $-\frac14F^a_{\mu\nu}F^{a\mu\nu}+J^{a\mu}A^a_\mu$ has first variation $\sum\delta A^a_\nu\,\mathcal E^{a\nu}-\partial_\mu(F^{a\mu\nu}\delta A^a_\nu)$, with $\mathcal E^{a\nu}=\nabla_\mu F^{a\mu\nu}+J^{a\nu}$.
\item If $A^2=A^3=0$ and $J^2=J^3=0$, the $a=1$ equation is Maxwell's equation.
\end{enumerate}
\end{theorem}
\begin{proof}
\emph{(i)} Differentiate $g\,g^{-1}=1$: $(\partial_kg)g^{-1}+g\,\partial_k(g^{-1})=0$, so $\partial_k(g^{-1})=-g^{-1}(\partial_kg)g^{-1}$.

For covariance, write $F_{ij}=\partial_iA_j-\partial_jA_i+[A_i,A_j]$ and $A'_i=gA_ig^{-1}+g\,\partial_ig^{-1}$. We use two consequences of the first identity:
\[
(\partial_ig^{-1})\,g=-g^{-1}\partial_ig,\qquad g\,\partial_ig^{-1}=-(\partial_ig)\,g^{-1}.
\]
\begin{itemize}[nosep]
\item Derivative terms. With $(i\leftrightarrow j)$ standing for the same expression with $i,j$ swapped,
\[
\partial_iA'_j-(i\leftrightarrow j)=\bigl[\partial_ig\,A_jg^{-1}+g\,\partial_iA_j\,g^{-1}+gA_j\partial_ig^{-1}+\partial_ig\,\partial_jg^{-1}\bigr]-(i\leftrightarrow j),
\]
because $g\,\partial_i\partial_jg^{-1}$ is symmetric and drops out.
\item Commutator terms. $[A'_i,A'_j]$ equals $g[A_i,A_j]g^{-1}$ plus
\[
\bigl[gA_i\partial_jg^{-1}-\partial_ig\,A_jg^{-1}-\partial_ig\,\partial_jg^{-1}\bigr]-(i\leftrightarrow j).
\]
\end{itemize}
Add the two lists.
\begin{itemize}[nosep]
\item Terms of the form $\partial g\,A\,g^{-1}$: from the derivative terms, $\partial_ig\,A_jg^{-1}-\partial_jg\,A_ig^{-1}$; from the commutator, $-\partial_ig\,A_jg^{-1}+\partial_jg\,A_ig^{-1}$. They cancel.
\item Terms of the form $g\,A\,\partial g^{-1}$: from the derivative terms, $gA_j\partial_ig^{-1}-gA_i\partial_jg^{-1}$; from the commutator, $gA_i\partial_jg^{-1}-gA_j\partial_ig^{-1}$. They cancel.
\item Terms of the form $\partial g\,\partial g^{-1}$: also cancel.
\end{itemize}
What remains is $g(\partial_iA_j-\partial_jA_i+[A_i,A_j])g^{-1}=gF_{ij}g^{-1}$. Only associativity was used, never commutativity.

\emph{(ii)} The derivative part is antisymmetric. For the quadratic part, rename $b\leftrightarrow c$ and use $\epsilon_{acb}=-\epsilon_{abc}$:
\[
\epsilon_{abc}A^b_\nu A^c_\mu=\epsilon_{acb}A^c_\nu A^b_\mu=-\epsilon_{abc}A^b_\mu A^c_\nu .
\]
For the variation, define the covariant derivative of an adjoint-valued one-form by $(\nabla_\mu X)^a_\nu=\partial_\mu X^a_\nu+\epsilon_{abc}A^b_\mu X^c_\nu$. The part of $F^a_{\mu\nu}(A+\delta A)$ that is linear in $\delta A$ is
\[
\partial_\mu\delta A^a_\nu-\partial_\nu\delta A^a_\mu+\epsilon_{abc}\bigl(\delta A^b_\mu A^c_\nu+A^b_\mu\delta A^c_\nu\bigr).
\]
Apply the same renaming to the first $\epsilon$ term: $\epsilon_{abc}\delta A^b_\mu A^c_\nu=-\epsilon_{abc}A^b_\nu\delta A^c_\mu$. The linear part is then exactly $(\nabla_\mu\delta A)^a_\nu-(\nabla_\nu\delta A)^a_\mu$.

\emph{(iii)} By (ii) and the antisymmetry of $F^{a\mu\nu}$,
\[
\delta\bigl(-\tfrac14F^2\bigr)=-\tfrac12F^{a\mu\nu}\bigl((\nabla_\mu\delta A)^a_\nu-(\nabla_\nu\delta A)^a_\mu\bigr)=-F^{a\mu\nu}(\nabla_\mu\delta A)^a_\nu .
\]
Expand the covariant derivative into two pieces.
\begin{itemize}[nosep]
\item Derivative piece: $-F^{a\mu\nu}\partial_\mu\delta A^a_\nu=(\partial_\mu F^{a\mu\nu})\delta A^a_\nu-\partial_\mu(F^{a\mu\nu}\delta A^a_\nu)$.
\item $\epsilon$ piece: $-\epsilon_{abc}F^{a\mu\nu}A^b_\mu\delta A^c_\nu$. Relabel the colour indices cyclically, $(a,b,c)\to(b,c,a)$. Since $\epsilon_{bca}=\epsilon_{abc}$, this is $-\epsilon_{abc}F^{b\mu\nu}A^c_\mu\delta A^a_\nu$. Swapping $b\leftrightarrow c$ turns it into $+\epsilon_{abc}A^b_\mu F^{c\mu\nu}\delta A^a_\nu$.
\end{itemize}
The coefficient of $\delta A^a_\nu$ is therefore $\partial_\mu F^{a\mu\nu}+\epsilon_{abc}A^b_\mu F^{c\mu\nu}+J^{a\nu}=\mathcal E^{a\nu}$.

\emph{(iv)} With only $A^1$ nonzero, every product $\epsilon_{abc}A^b A^c$ needs $b=c=1$ and vanishes. Hence $F^2=F^3=0$ and $F^1=\dd A^1$. In $\mathcal E^{1\nu}$ the extra term $\epsilon_{1bc}A^bF^c$ needs $b=1$ and vanishes. In $\mathcal E^{2\nu}$ and $\mathcal E^{3\nu}$ every term vanishes. So the only equation left is $\partial_\mu F^{1\mu\nu}+J^{1\nu}=0$, which is Maxwell's equation (\cref{thm:maxwellvar}).
\end{proof}


The bookkeeping lemmas used in (iii), namely the relabelling of $\epsilon_{abc}$ in a triple product and $\sum F^{\mu\nu}(X_{\mu\nu}-X_{\nu\mu})=2\sum F^{\mu\nu}X_{\mu\nu}$, are recorded in the checklist.

\subsection*{Homogeneous fields on the three-sphere}

On $S^3=\SU(2)$ the simplest Yang--Mills fields are multiples of the Maurer--Cartan form, $A=\lambda\,U^{-1}\dd U$. They are invariant under left translations, so the Yang--Mills equation reduces to an algebraic condition on $\lambda$.

\begin{theorem}[Homogeneous Yang--Mills fields on $S^3$]\label{thm:merons}\ver\
For $A=\lambda\,U^{-1}\dd U$ the Yang--Mills expression is $8\ii\lambda(1-\lambda)(2\lambda-1)$ times a fixed generator. Hence $A$ solves the Yang--Mills equation if and only if $\lambda\in\{0,\tfrac12,1\}$.
\end{theorem}
\begin{proof}
The computation is done in a left-invariant frame, so no coordinates on $S^3$ are needed. It has four steps.

\emph{Step 1: frame and connection.} Let $s_a=\sigma_a/2$, so that $[s_a,s_b]=\ii\epsilon_{abc}s_c$. The left-invariant vector fields $X_a$ that correspond to $\ii\sigma_a=2\ii s_a$ satisfy
\[
[X_a,X_b]=-2\epsilon_{abc}X_c ,
\]
because $[\ii\sigma_a,\ii\sigma_b]=-2\ii\epsilon_{abc}\sigma_c$. Normalise the bi-invariant metric so that the $X_a$ are orthonormal. The Levi-Civita connection of a bi-invariant metric is $\nabla_XY=\frac12[X,Y]$ on left-invariant fields.

The Maurer--Cartan form sends $X_a$ to its Lie algebra element, $U^{-1}\dd U(X_a)=2\ii s_a$. Therefore the frame components of $A$ are the constants $A_a=2\ii\lambda\,s_a$.

\emph{Step 2: curvature.} The components $A_a$ are constant, so $X_a(A_b)=0$. The definition $F(X,Y)=X(A(Y))-Y(A(X))-A([X,Y])+[A(X),A(Y)]$ then gives
\[
F_{ab}=[A_a,A_b]+2\epsilon_{abc}A_c=(2\ii\lambda)^2\ii\epsilon_{abc}s_c+4\ii\lambda\epsilon_{abc}s_c=4\ii\lambda(1-\lambda)\,\epsilon_{abc}s_c .
\]
This vanishes at $\lambda=0,1$, as it must for a pure gauge.

\emph{Step 3: divergence of a left-invariant two-form.} The components $F_{ab}$ are constant, so
\[
(\nabla_{X_a}F)(X_a,X_b)=-F(\nabla_{X_a}X_a,X_b)-F(X_a,\nabla_{X_a}X_b)=-F\bigl(X_a,-\epsilon_{abc}X_c\bigr)=\epsilon_{abc}F_{ac},
\]
since $\nabla_{X_a}X_a=\frac12[X_a,X_a]=0$. The Yang--Mills expression in direction $b$ is therefore
\[
\mathrm{YM}_b=\sum_a\Bigl([A_a,F_{ab}]+\sum_c\epsilon_{abc}F_{ac}\Bigr).
\]

\emph{Step 4: evaluate.} We need the contraction
\[
\sum_{a,c}\epsilon_{abc}\epsilon_{acd}=-\sum_{a,c}\epsilon_{abc}\epsilon_{adc}=-2\delta_{bd}.
\]
Then:
\begin{align*}
\sum_a[A_a,F_{ab}]&=2\ii\lambda\cdot4\ii\lambda(1-\lambda)\sum_{a,c}\epsilon_{abc}[s_a,s_c]
=-8\lambda^2(1-\lambda)\,\ii\sum_{a,c,d}\epsilon_{abc}\epsilon_{acd}s_d=16\ii\lambda^2(1-\lambda)\,s_b,\\
\sum_{a,c}\epsilon_{abc}F_{ac}&=4\ii\lambda(1-\lambda)\sum_{a,c,d}\epsilon_{abc}\epsilon_{acd}s_d=-8\ii\lambda(1-\lambda)\,s_b .
\end{align*}
The sum is $8\ii\lambda(1-\lambda)(2\lambda-1)\,s_b$. The matrix $s_b$ is nonzero, so $\mathrm{YM}_b=0$ for all $b$ if and only if $\lambda(1-\lambda)(2\lambda-1)=0$, that is, $\lambda\in\{0,\frac12,1\}$.
\end{proof}


The values $\lambda=0$ and $\lambda=1$ are the trivial connection and its gauge transform, both flat. The value $\lambda=\frac12$ is the meron \cite{deAlfaroFubiniFurlan1976}, half of a pure gauge. Its field strength does not vanish, and it will reappear in Chapter~15 as the matter content of a closed universe.

\section{The spherically symmetric magnetic ansatz}

The general spherically symmetric $\SU(2)$ potential can be brought, in a suitable gauge and in the purely magnetic sector, to the form \cite{Witten1977,BartnikMcKinnon1988}
\begin{equation}\label{eq:magansatz}
A^1=w\,\dd\theta,\qquad A^2=w\sin\theta\,\dd\phi,\qquad A^3=\cos\theta\,\dd\phi,
\end{equation}
with a single function $w(t,x)$. The value $w=\pm1$ is a pure gauge (vacuum), and $w=0$ is the embedded Dirac monopole of unit charge (\cref{thm:dirac}).

\begin{theorem}[Field strength and stress tensor of the magnetic ansatz]\label{thm:ymstress}\ver\
\begin{enumerate}[label=(\roman*)]
\item The field strength of \eqref{eq:magansatz} has components $F^1_{t\theta}=w_t$, $F^1_{x\theta}=w_x$, $F^2_{t\phi}=w_t\sin\theta$, $F^2_{x\phi}=w_x\sin\theta$ and $F^3_{\theta\phi}=(w^2-1)\sin\theta$, and no others.
\item With $\hat T_{\mu\nu}=F^a_{\mu\lambda}F^{a}{}_\nu{}^\lambda-\frac14g_{\mu\nu}F^a_{\alpha\beta}F^{a\alpha\beta}$ and the diagonal metric \eqref{eq:diagmetric}, the components $\hat T_{tt}$, $\hat T_{xx}$, $\hat T_{tx}=2w_tw_x/R^2$ and $\hat T_{\theta\theta}=(1-w^2)^2/(2R^2)$ are explicit; in particular
\[
\frac{\hat T_{tt}}{N^2}=\frac{w_t^2}{N^2R^2}+\frac{w_x^2}{A^2R^2}+\frac{(1-w^2)^2}{2R^4}.
\]
\item In the double-null gauge \eqref{eq:nullmetric}, $\hat T_{uu}=2w_u^2/r^2$, $\hat T_{vv}=2w_v^2/r^2$, $\hat T_{uv}=\Omega(1-w^2)^2/(2r^4)$, $\hat T_{\theta\theta}=(1-w^2)^2/(2r^2)$, and the mixed angular components vanish.
\end{enumerate}
\end{theorem}
\begin{proof}
\emph{(i)} Use $F^a_{\mu\nu}=\partial_\mu A^a_\nu-\partial_\nu A^a_\mu+\epsilon_{abc}A^b_\mu A^c_\nu$, the coordinates $(t,x,\theta,\phi)$ and $w=w(t,x)$. The only nonzero potential components are $A^1_\theta=w$, $A^2_\phi=w\sin\theta$ and $A^3_\phi=\cos\theta$. Since $A_t=A_x=0$, the quadratic term can be nonzero only for the pair $(\theta,\phi)$.
\begin{itemize}[nosep]
\item $F^1_{t\theta}=\partial_tw$ and $F^1_{x\theta}=\partial_xw$.
\item $F^2_{t\phi}=w_t\sin\theta$ and $F^2_{x\phi}=w_x\sin\theta$.
\item $F^1_{\theta\phi}=-\partial_\phi w+\epsilon_{123}(A^2_\theta A^3_\phi-A^3_\theta A^2_\phi)=0$, since $A^2_\theta=A^3_\theta=0$.
\item $F^2_{\theta\phi}=\partial_\theta(w\sin\theta)+\epsilon_{231}A^3_\theta A^1_\phi+\epsilon_{213}A^1_\theta A^3_\phi=w\cos\theta+0-w\cos\theta=0$.
\item $F^3_{\theta\phi}=\partial_\theta\cos\theta+\epsilon_{312}A^1_\theta A^2_\phi=-\sin\theta+w^2\sin\theta$.
\end{itemize}
All other components vanish because $A^3$ does not depend on $t$ or $x$.

\emph{(ii)} The inverse metric is diagonal: $g^{tt}=-1/N^2$, $g^{xx}=1/A^2$, $g^{\theta\theta}=1/R^2$, $g^{\phi\phi}=1/(R^2\sin^2\theta)$. Then
\[
F^a_{\alpha\beta}F^{a\alpha\beta}=2\sum_{\alpha<\beta}g^{\alpha\alpha}g^{\beta\beta}(F^a_{\alpha\beta})^2
=2\Bigl(-\frac{2w_t^2}{N^2R^2}+\frac{2w_x^2}{A^2R^2}+\frac{(1-w^2)^2}{R^4}\Bigr).
\]
The $\sin^2\theta$ from $F^2$ and $F^3$ cancels against $g^{\phi\phi}$. Next,
\[
F^a_{t\lambda}F^{a}{}_t{}^\lambda=\frac{w_t^2}{R^2}+\frac{w_t^2\sin^2\theta}{R^2\sin^2\theta}=\frac{2w_t^2}{R^2},
\]
and $-\frac14g_{tt}F^2=\frac{N^2}4F^2$. Adding,
\[
\hat T_{tt}=\frac{2w_t^2}{R^2}-\frac{w_t^2}{R^2}+\frac{N^2w_x^2}{A^2R^2}+\frac{N^2(1-w^2)^2}{2R^4},
\]
and dividing by $N^2$ gives the displayed formula. The same steps give
\[
\hat T_{xx}=\frac{w_x^2}{R^2}+\frac{A^2w_t^2}{N^2R^2}-\frac{A^2(1-w^2)^2}{2R^4}.
\]
The off-diagonal component is $\hat T_{tx}=g^{\theta\theta}F^1_{t\theta}F^1_{x\theta}+g^{\phi\phi}F^2_{t\phi}F^2_{x\phi}=2w_tw_x/R^2$. For the angular component,
\[
\hat T_{\theta\theta}=\Bigl(-\frac{w_t^2}{N^2}+\frac{w_x^2}{A^2}+\frac{(1-w^2)^2}{R^2}\Bigr)-\frac{R^2}4F^2=\frac{(1-w^2)^2}{2R^2}.
\]
Here the kinetic and gradient terms cancel exactly, and half of the magnetic term survives.

\emph{(iii)} Now $g_{uv}=-\Omega$ and $g^{uv}=-1/\Omega$, and the components of (i) hold with $(t,x)$ replaced by $(u,v)$. Because $g^{uu}=g^{vv}=0$, every contraction over a null index pairs $u$ with $v$:
\[
F^a_{\alpha\beta}F^{a\alpha\beta}=4g^{uv}g^{\theta\theta}F^1_{u\theta}F^1_{v\theta}+4g^{uv}g^{\phi\phi}F^2_{u\phi}F^2_{v\phi}+\frac{2(1-w^2)^2}{r^4}=-\frac{8w_uw_v}{\Omega r^2}+\frac{2(1-w^2)^2}{r^4}.
\]
Then the components follow one by one.
\begin{itemize}[nosep]
\item $g_{uu}=0$, so $\hat T_{uu}=g^{\theta\theta}(F^1_{u\theta})^2+g^{\phi\phi}(F^2_{u\phi})^2=2w_u^2/r^2$. The same holds for $\hat T_{vv}$.
\item $\hat T_{uv}=\frac{2w_uw_v}{r^2}+\frac\Omega4F^2=\frac{\Omega(1-w^2)^2}{2r^4}$.
\item $\hat T_{\theta\theta}=2g^{uv}F^1_{\theta u}F^1_{\theta v}+\frac{(1-w^2)^2}{r^2}-\frac{r^2}4F^2=\frac{(1-w^2)^2}{2r^2}$.
\end{itemize}
For $\hat T_{\theta\phi}$ and $\hat T_{u\theta}$, every term has a factor that is zero in the list (i). For example, $F^1_{u\lambda}$ is nonzero only for $\lambda=\theta$, but $F^1_{\theta\theta}=0$. These components therefore vanish.
\end{proof}


The physical stress tensor is $T_{\mu\nu}=\hat T_{\mu\nu}/g^2$. The energy density is the sum of a kinetic term, a gradient term and the magnetic self-energy $(1-w^2)^2/(2g^2R^4)$ of the angular field. The last term is what distinguishes Yang--Mills from Maxwell in spherical symmetry: at $w=0$ it is the Coulomb energy of a unit magnetic charge, and at $w=\pm1$ it vanishes.

\begin{litbox}
All equations in this chapter are standard; see Jackson for Maxwell \cite{Jackson}, Landau and Lifshitz for fluids \cite{LandauFluid}, and for Yang--Mills fields in spherical symmetry the ansatz of Witten \cite{Witten1977} as used by Bartnik and McKinnon \cite{BartnikMcKinnon1988} and reviewed by Volkov and Gal'tsov \cite{VolkovGaltsov1999}. The meron is due to de Alfaro, Fubini and Furlan \cite{deAlfaroFubiniFurlan1976}. The contribution here is organisational: each equation is labelled as an off-shell identity or an on-shell consequence, and the stress tensor of the magnetic ansatz is derived in both gauges that Part~IV uses.
\end{litbox}

\section*{What this chapter passes on}
\begin{itemize}[nosep]
\item Maxwell and Yang--Mills equations as Euler--Lagrange equations; Bianchi identities; conservation laws (\cref{thm:bianchi,thm:maxwellvar,thm:ymvar}).
\item The hemisphere point of a classical charged particle (\cref{thm:particle}).
\item Homogeneous Yang--Mills fields on $S^3$: $\lambda\in\{0,\frac12,1\}$ (\cref{thm:merons}).
\item The magnetic ansatz and its stress tensor (\cref{thm:ymstress}).
\end{itemize}

\section*{Exercises}
\begin{exercise}
From $F_{\mu\nu}=\partial_\mu A_\nu-\partial_\nu A_\mu$ derive $\nabla\cdot\mathbf B=0$ and $\nabla\times\mathbf E+\partial_t\mathbf B=0$ with $E_i=F_{i0}$, $B_i=\frac12\epsilon_{ijk}F_{jk}$ (up to the sign conventions of your choice), and compare with \cref{thm:bianchi}.
\end{exercise}
\begin{exercise}
Compute $F^3_{\theta\phi}$ for \eqref{eq:magansatz} and check $(w^2-1)\sin\theta$.
\end{exercise}
\begin{exercise}
Show that $\lambda(1-\lambda)(2\lambda-1)=0$ has exactly the three roots of \cref{thm:merons}.
\end{exercise}
\begin{exercise}
Verify $[v,w]=2v\times w$ for $v=i$, $w=j$.
\end{exercise}

\begin{checklist}
\chk{AlphaVII_Maxwell:bianchi}{Bianchi identity $\partial_{[\lambda}F_{\mu\nu]}=0$.$^*$}
\chk{AlphaIX_Classical:bianchi_div}{$\nabla\cdot\mathbf B=0$ identically.$^*$}
\chk{AlphaIX_Classical:faraday}{$\nabla\times\mathbf E+\partial_t\mathbf B=0$ identically.$^*$}
\chk{AlphaVII_Maxwell:div_div_F}{$\partial_\nu\partial_\mu F^{\mu\nu}=0$.$^*$}
\chk{AlphaIX_Classical:eta_const,AlphaIX_Classical:eta_const',AlphaIX_Classical:d_quarter,AlphaIX_Fluid:d_eta}{The metric components and numerical constants have zero derivative (four statements).}
\chk{AlphaVII_Maxwell:first_variation}{First variation of the Maxwell Lagrangian, exact to second order.$^*$}
\chk{AlphaIX_Classical:first_variation}{First variation in Euler--Lagrange plus divergence form.$^*$}
\chk{AlphaIX_Classical:gauss}{$\mathcal E^0=0\iff\nabla\cdot\mathbf E=\rho$.$^*$}
\chk{AlphaIX_Classical:ampere}{$\mathcal E^1=\partial_tE_1-(\nabla\times\mathbf B)_1+J^1$.$^*$}
\chk{AlphaVII_Maxwell:current_conserved}{On shell, $\partial_\nu J^\nu=0$.$^*$}
\chk{AlphaIX_Classical:charge_conservation}{On shell, $\partial_\nu J^\nu=0$ (second formulation).$^*$}
\chk{AlphaIX_Classical:gauss_integral}{Integrated spherical Gauss law.$^*$}
\chk{AlphaIX_Classical:stress_divergence}{$\partial^\mu T_{\mu\nu}=F_{\nu\alpha}\partial_\mu F^{\mu\alpha}$.$^*$}
\chk{AlphaIX_Classical:particle_el}{Euler--Lagrange equations of a charged particle: the Lorentz force.$^*$}
\chk{AlphaIX_Classical:hamilton}{Hamilton's equations are equivalent to the Lorentz force law.$^*$}
\chk{AlphaIX_Classical:mass_shell}{Kinetic momentum on the mass shell, and the hemisphere point $(m/E,\mathbf p/E)$.$^*$}
\chk{AlphaIX_Classical:dust_lorentz}{Charged dust obeys the Lorentz force law.$^*$}
\chk{AlphaIX_Fluid:divT_expand}{Expansion of $\partial_\mu T^{\mu\nu}$ for a perfect fluid.$^*$}
\chk{AlphaIX_Fluid:norm_deriv}{$u\cdot u=-1$ implies $u_\nu\partial_\mu u^\nu=0$.$^*$}
\chk{AlphaIX_Fluid:fluid_energy}{Energy equation $\partial_\mu((\rho+p)u^\mu)=u^\mu\partial_\mu p$.$^*$}
\chk{AlphaIX_Fluid:fluid_euler}{Relativistic Euler equation.$^*$}
\chk{AlphaIX_Fluid:cyclic_zero}{$\dd\omega=0$ in components.$^*$}
\chk{AlphaIX_Fluid:lie_two_form}{Cartan's formula $\mathcal L_u\omega=\dd(i_u\omega)$ for closed $\omega$.$^*$}
\chk{AlphaIX_Fluid:kelvin}{Kelvin: $i_u\omega=0\Rightarrow\mathcal L_u\omega=0$.$^*$}
\chk{AlphaIX_Fluid:flux_freezing}{Flux freezing in ideal MHD.$^*$}
\chk{AlphaIX_Fluid:cross_is_commutator}{$[v,w]=2v\times w$ for pure quaternions.$^*$}
\chk{AlphaIX_Fluid:helicity_balance}{Magnetic helicity balance.$^*$}
\chk{GaugeCovariance:d_inv}{$\partial_k(g^{-1})=-g^{-1}(\partial_kg)g^{-1}$.$^*$}
\chk{GaugeCovariance:curvature_covariant}{$F'=gFg^{-1}$ under gauge transformations.$^*$}
\chk{AlphaIX_YM:Fa_antisym,AlphaIX_YM:Fau_antisym,AlphaIX_YM:Fau_diag}{Antisymmetry of $F^a_{\mu\nu}$ with lower and upper indices, and vanishing diagonal (three statements).$^*$}
\chk{AlphaIX_YM:dF_cov}{$\delta F=\nabla\delta A-\nabla\delta A$ (antisymmetrised).$^*$}
\chk{AlphaIX_YM:antisym_sum}{$\sum F^{\mu\nu}(X_{\mu\nu}-X_{\nu\mu})=2\sum F^{\mu\nu}X_{\mu\nu}$.}
\chk{AlphaIX_YM:relabel}{Relabelling of $\epsilon_{abc}$ in triple products.}
\chk{AlphaIX_YM:ym_first_variation}{First variation of the Yang--Mills Lagrangian.$^*$}
\chk{AlphaIX_YM:u1_projection}{Projection to $U(1)\subset\SU(2)$ gives Maxwell's equation.$^*$}
\chk{AlphaVII_Bundle:ym_formula}{Yang--Mills expression for $\lambda U^{-1}\dd U$ on $S^3$: $8\ii\lambda(1-\lambda)(2\lambda-1)$ times a generator.$^*$}
\chk{AlphaVII_Bundle:ym_roots}{Solutions exactly for $\lambda\in\{0,\frac12,1\}$.$^*$}
\chk{AlphaS4_Curv:fym_eq}{Field strength of the magnetic ansatz.$^*$}
\chk{AlphaS4_Curv:tymD_00,AlphaS4_Curv:tymD_11,AlphaS4_Curv:tymD_01,AlphaS4_Curv:tymD_22}{Stress tensor of the magnetic ansatz in the diagonal gauge (four components).$^*$}
\chk{AlphaS4_Curv:tymN_00,AlphaS4_Curv:tymN_11,AlphaS4_Curv:tymN_01,AlphaS4_Curv:tymN_22,AlphaS4_Curv:tymN_12,AlphaS4_Curv:tymN_13}{Stress tensor of the magnetic ansatz in the double-null gauge (six components).$^*$}
\end{checklist}
